Nibli integrates a formal predicate-call language, bounded symbolic evaluation,
selective materialization, retractable premise identities, and optional evidence
envelopes across native and WebAssembly execution. We describe its supported
reasoning model and distinguish derivation from accepted premises from claims
about input truth. Selected Lean proofs establish conditional model properties;
implementation tests connect those models to code, while envelope validation
checks structural coherence. A frozen research artifact compares finite Horn
and stratified-negation workloads with clingo and compiled, single-threaded
\Souffle, varies evaluation bounds, and exercises withdrawal and saved-state
reopening. Across \MeasuredRuns\ measured processes, no definitive-answer
disagreement is observed, while \TimeoutRuns\ executions reach an external
cutoff. The audit checks \CertificateCount\ envelopes and \CitationCount\ source
citations without a coherence or source-liveness error. Materialization recovers
\BoundsGains\ returned definitive answers in the fixed bounds slice. At the
1,000-fact target, certificate query time is \CertificateMedianRatio\ times
verdict-only query time by ratio of medians. Larger ingestion and forward
recursion cases expose substantial limits under the cutoff. Native, Wasmtime,
and Node/V8 agree on all \PortableAnswers\ measured query outcomes in the shared
scenarios. The results characterize an integration whose auditability depends
on explicit premise identity, evaluation profiles, and evidence boundaries,
while its performance and completeness remain workload-dependent.
